% Ban thu: chi giu cac cau co trong CD2

%@ id=TT25.thucchien01.003 | trang=1 | so=II | de=II | vitri=CD2.5,C3.9.6 | muc=9 | tt=0
\begin{bt}
Cho hai biểu thức: $A = \dfrac{\sqrt{x}+2}{\sqrt{x}-3}$ và $B = \dfrac{\sqrt{x}-2}{\sqrt{x}-3} - \dfrac{3}{\sqrt{x}} + \dfrac{x+4\sqrt{x}-9}{x-3\sqrt{x}}$ (với $x > 0; x \ne 9$).
\begin{cacy}
\item Tính giá trị của biểu thức $A$ khi $x = 169$.
\item Chứng minh $B = \dfrac{2\sqrt{x}-1}{\sqrt{x}-3}$.
\item Tìm giá trị nguyên của $x$ để biểu thức $P = B : A$ đạt giá trị nhỏ nhất.
\end{cacy}
\loigiai{
\begin{cacy}
\item Thay $x = 169$ (thoả mãn ĐKXĐ) vào biểu thức $A$ ta được:
\[ A = \dfrac{15}{10} = \dfrac{3}{2}. \]
Kết luận.
\item Ta có:
\[ B = \dfrac{\sqrt{x}-2}{\sqrt{x}-3} - \dfrac{3}{\sqrt{x}} + \dfrac{x+4\sqrt{x}-9}{\sqrt{x}(\sqrt{x}-3)} \]
\[ B = \dfrac{x-2\sqrt{x}-3\sqrt{x}+9+x+4\sqrt{x}-9}{\sqrt{x}(\sqrt{x}-3)} \]
\[ B = \dfrac{2x-\sqrt{x}}{\sqrt{x}(\sqrt{x}-3)} = \dfrac{\sqrt{x}(2\sqrt{x}-1)}{\sqrt{x}(\sqrt{x}-3)} \]
\[ B = \dfrac{2\sqrt{x}-1}{\sqrt{x}-3}. \]
Vậy với $x > 0; x \ne 9$ thì $B = \dfrac{2\sqrt{x}-1}{\sqrt{x}-3}$.
\item Với $x > 0; x \ne 9$ ta có:
\[ P = B : A = \dfrac{2\sqrt{x}-1}{\sqrt{x}-3} \cdot \dfrac{\sqrt{x}-3}{\sqrt{x}+2} = \dfrac{2\sqrt{x}-1}{\sqrt{x}+2} = 2 - \dfrac{5}{\sqrt{x}+2}. \]
Vì $x > 0; x \ne 9$ và $x$ nguyên nên $x \ge 1$. Suy ra $\sqrt{x}+2 \ge 3$.

Lập luận: $P \ge 2 - \dfrac{5}{3} = \dfrac{1}{3}$.

Dấu “=” xảy ra khi $x = 1$ (TMĐK).

Vậy $\min P = \dfrac{1}{3}$ tại $x = 1$.
\end{cacy}
}
\end{bt}
